\documentclass[10pt]{article}
\usepackage[T1]{fontenc}
\usepackage{lmodern}
\usepackage[margin=0.75in]{geometry}
\usepackage{amsmath,amssymb,amsthm,booktabs,microtype}
\usepackage[hidelinks]{hyperref}
\setlength{\parindent}{0pt}
\setlength{\parskip}{3pt}
\newcommand{\Q}{\mathbb Q}
\newcommand{\HH}{\operatorname{HH}}
\newcommand{\End}{\operatorname{End}}
\newcommand{\Hom}{\operatorname{Hom}}
\newtheorem{theorem}{Theorem}
\title{Vanishing first Hochschild cohomology with infinite representation type\\\large A counterexample to Problem 00000000556}
\author{lizaixi01}
\date{October 24, 2026}
\begin{document}
\maketitle
\vspace{-1.5em}

\textbf{Scope.} The printed biconditional implies
\[
 \HH^1(\Lambda)=0\quad\Longrightarrow\quad
 \Lambda\text{ is representation-finite}.
\]
This necessary implication is independent of the meaning of the additional
adjective ``rigid.'' We disprove it for a finite-dimensional algebra over $\Q$.
Representation-finite means that the finitely generated indecomposable left
modules have only finitely many isomorphism classes. The source's parenthetical
mentions verification for string algebras. Our four-leaf star has four arrows
entering its central vertex, violating the at-most-two-incoming condition for
string algebras. The conclusion concerns the printed general algebra statement;
it does not refute a statement restricted to that subclass.

\begin{theorem}
There is a nine-dimensional rational path algebra $A$ with $\HH^1(A)=0$ and
infinitely many pairwise nonisomorphic indecomposable six-dimensional left
$A$-modules. Consequently the printed biconditional fails for every possible
additional rigidity predicate.
\end{theorem}

\textbf{The actual algebra.}
Take the five vertices $0,1,2,3,4$ and four arrows $a_j:j+1\to0$,
$j=0,1,2,3$. This is the four-leaf extended Dynkin quiver $\widetilde D_4$.
There are five length-zero paths $e_i$, four length-one paths $a_j$, and no
longer paths. The path-algebra definition uses all paths as a basis and actual
concatenation as multiplication~\cite[p.~168]{EHIS}.

Lean constructs $A$ as the native square-zero extension
$\Q^5\ltimes\Q^4$, with faithful multiplication
\[
 (d,b)(d',b')=
 \left((d_i d'_i)_{i=0}^4,
       (d_0 b'_j+b_jd'_{j+1})_{j=0}^3\right).
 \tag{1}
\]
It enumerates \emph{all} native quiver paths, proves their lengths and uniqueness,
identifies their basis with the basis of this extension, and proves both
concatenation multiplication and zero multiplication for noncomposable paths.
Thus (1) is the actual full path algebra, not an assigned dimension model.
In particular
\[
 e_0a_j=a_j=a_je_{j+1},\qquad a_ja_k=0,\qquad 1=\sum_{i=0}^4e_i,
 \qquad\dim_{\Q}A=9.
\]

\textbf{The true first Hochschild cohomology vanishes.}
Use the usual regular $A$-bimodule coefficients~\cite[\S1.5, pp.~37--38]{Loday}:
\[
 C^0=A,\quad C^1=\Hom_{\Q}(A,A),\quad
 C^2=\Hom_{\Q}(A\otimes_{\Q}A,A),
\]
\[
 (d^0h)(x)=xh-hx,\qquad
 (d^1f)(x,y)=xf(y)-f(xy)+f(x)y.
 \tag{2}
\]
The formalization has the actual tensor--Hom equivalence with curried bilinear
maps, proves the displayed formula on every tensor $x\otimes y$, and verifies
$d^1d^0=0$. It forms the native short complex in $\texttt{ModuleCat}\,\Q$;
$\HH^1(A)$ is its actual homology object.

If $D$ is a 1-cocycle, (2) says exactly that $D$ is a rational derivation.
Applying the Leibniz rule to $e_i^2=e_i$ makes every diagonal coordinate of
$D(e_i)$ zero. The relations $e_0e_{j+1}=0$ and
$e_0a_j=a_j=a_je_{j+1}$ then give, for coefficients obtained from the actual $D$,
\[
 D(e_0)=\sum_j u_ja_j,\qquad D(e_{j+1})=-u_ja_j,
 \qquad D(a_j)=\lambda_ja_j.
\]
Put $h=\sum_j\lambda_je_{j+1}+\sum_j u_ja_j$. Formula (1) gives
$D(e_i)=e_ih-he_i$ and $D(a_j)=a_jh-ha_j$. Agreement on the full path basis
proves $D=d^0h$. Every actual cocycle is therefore a boundary. Lean proves
exactness of the native complex and hence $\HH^1(A)=0$.

\newpage
\textbf{A family of genuine finite-dimensional modules.}
For $t\in\Q$, let $M_t$ have central space $\Q^2$ and a copy of $\Q$ at each
of the four leaves. The four arrow maps send $1$ to
\[
 v_0=(1,0),\quad v_1=(0,1),\quad v_2=(1,1),\quad v_3(t)=(1,t).
 \tag{3}
\]
More explicitly, the actual carrier is $\Q^2\oplus\Q^4$. Writing a vector as
$(z,u)$, with $z\in\Q^2$ and $u=(u_j)$, define
\[
 (d,b)\cdot(z,u)=
 \left(d_0z+\sum_j b_ju_jv_j(t),\ (d_{j+1}u_j)_j\right).
 \tag{4}
\]
Equation (1) directly proves the unit, associativity, addition and scalar laws
for (4). Lean constructs the native left $A$-module instance, its compatible
rational scalar structure, and its actual dimension $6$. It also proves finite
generation as an $A$-module.

\textbf{All native morphisms are controlled.}
Let $f:M_t\to M_s$ be an arbitrary $A$-linear map. Commutation with every
actual vertex idempotent makes it a central linear map $T:\Q^2\to\Q^2$
together with four leaf scalars $c_j$. Commutation with the actual arrows says
\[
 T v_j(t)=c_jv_j(s)\qquad (j=0,1,2,3).
 \tag{5}
\]
The first two equations make $T=\operatorname{diag}(c_0,c_1)$.
The third gives $c_0=c_1=c_2$. The first coordinate of the fourth gives
$c_3=c_0$, and its second coordinate gives $c_0t=c_0s$.
Thus, on all six actual coordinates, $f$ is the scalar map with one common
scalar $c$, satisfying
\[
 c(t-s)=0.
 \tag{6}
\]
These are proved for every native module morphism, using all module basis
vectors, not just a chosen class of quiver maps.

For $t\ne s$, (6) forces every $A$-linear map $M_t\to M_s$ to be zero.
Since $M_t\ne0$, the modules cannot be isomorphic. For $s=t$, every
endomorphism of $M_t$ is a rational scalar.

To prove actual indecomposability, let $P,Q$ be any complementary native
$A$-submodules of $M_t$. Their actual projection $p:M_t\to M_t$ is scalar,
say $p=c\,\mathrm{id}$. If $P\ne0$, take a nonzero vector of $P$; its projection
is itself, so $c=1$. Projection kills every vector of $Q$, so $Q=0$.
This proves indecomposability against every direct-sum decomposition.

\textbf{The full isomorphism-class obstruction.}
The formal class space contains all finitely generated indecomposable objects
of the native category $\texttt{ModuleCat}\,A$, in a fixed small universe,
quotiented by the existence of a native $A$-linear equivalence. Finite generation
over the finite-dimensional algebra $A$ is proved equivalent to finite
dimension over $\Q$. The actual objects $M_t$ define an injection
\[
 \Q\lhook\joinrel\longrightarrow
 \{\text{finite-dimensional indecomposable left }A\text{-modules}\}/{\cong}.
\]
Hence this full class space is infinite and $A$ is not representation-finite.
The obstruction also holds in any larger universe, since it already occurs
in the displayed family.

\textbf{The printed assertion fails independently of rigidity.}
For any proposition $R$ assigned to this particular algebra, the proposed
biconditional
\[
 \HH^1(A)=0\quad\Longleftrightarrow\quad
 \bigl(A\text{ is representation-finite}\bigr)\land R
\]
would imply representation-finiteness from the true left side, contradicting
the injection above. The Lean theorem \texttt{printed\_iff\_impossible}
quantifies over every such proposition. No numerical invariant, cohomology
group, representation count or meaning of rigidity is stipulated to force
the conclusion.

\newpage
\textbf{The error in the previous submission.}
\href{https://github.com/The-Last-Math-Competition/The-Last-Math-Competition/pull/134}{PR~134}
attempted the opposite implication with the dual numbers $k[x]/(x^2)$.
Its non-inner derivation is a valid warning about Hochschild cohomology, but
it replaced ``rigid'' by Frobenius/self-injective. Those properties do not
establish the rigidity premise in the printed assertion. The reviewer
\href{https://github.com/The-Last-Math-Competition/The-Last-Math-Competition/pull/134#issuecomment-5942503126}{rejected that substitution}:
the dual numbers have a two-dimensional center and nonzero second Hochschild
cohomology in characteristic zero, so the proposed example does not meet
the center or deformation interpretations discussed there. The author
\href{https://github.com/The-Last-Math-Competition/The-Last-Math-Competition/pull/134#issuecomment-5945753401}{accepted the correction and withdrew the example}.
The PR is closed and unmerged.

The previous Lean file used $\mathbb N\times\mathbb N$ arithmetic rather than
a finite-dimensional algebra over a field. It did not construct the required
field-based Hochschild cochain homology or a classification of the full
finite module category establishing its representation-finiteness claim.
That certificate therefore did not repair the missing premise. The present
proof uses a rational algebra, the true Hochschild complex, actual module
instances and the full native isomorphism-class quotient. It attacks the
necessary implication in the other direction and works without selecting
any interpretation of ``rigid.'' It does not reuse the withdrawn argument.

\textbf{Formalization and the definition boundary.}
All of the following are built and proved in Lean: the actual square-zero
algebra and rational dimension; every native path and its basis and
multiplication; the full linear cochains in degrees $0,1,2$, the standard
bar coboundaries and tensor formula; exactness and native first homology;
the genuine family modules, finite dimensions and all native module
morphisms; arbitrary complementary-submodule indecomposability; and the
injection into the quotient of all finite indecomposable module objects.

The identification of first Hochschild cohomology with the degree-one
homology of (2) is the standard defining cochain model in
Loday~\cite[Definition~1.5.1 and \S1.5.2]{Loday}. Only the degrees needed for
that homology are constructed. No higher Hochschild group or general
derived-functor comparison theorem is asserted. The complete bilingual
conjecture is preserved in \texttt{SOURCE.md}.

\textbf{Reproduction and evidence.}
The project under \texttt{lean/} pins Lean~4.33.0 and Mathlib commit:
\begin{center}
{\footnotesize\texttt{db584cd6d46c92f209a44c0f1c829460d327499d}}
\end{center}
Run the four commands separately:
\begin{verbatim}
lean --version
lake build +Main
lake env lean -DwarningAsError=true Check.lean
lake env leanchecker --fresh Main
\end{verbatim}
The dynamic audit emits ordinary \texttt{\#print axioms} reports for every
actual declaration in \texttt{TLMC556}, including generated names. The
command logs, source bindings, execution receipts, actual thread-mask
samples and SHA-256 manifest accompany the submission. Only
\texttt{propext}, \texttt{Classical.choice} and \texttt{Quot.sound} are permitted.
Production timings are diagnostic and imply no controlled efficiency gain.
Organizer review remains pending.

\begin{thebibliography}{2}\small
\bibitem{Loday} J.-L.~Loday, \emph{Cyclic Homology}, second edition,
Grundlehren der mathematischen Wissenschaften~301, Springer, 1998,
pp.~37--38, Definition~1.5.1 and \S1.5.2.
\href{https://doi.org/10.1007/978-3-662-11389-9}{DOI: 10.1007/978-3-662-11389-9}.
\bibitem{EHIS} K.~Erdmann, T.~Holm, O.~Iyama and J.~Schr\"oer,
\emph{Radical embeddings and representation dimension},
Advances in Mathematics~185 (2004), 159--177, p.~168 (path-algebra definition).
\href{https://doi.org/10.1016/S0001-8708(03)00169-5}{DOI: 10.1016/S0001-8708(03)00169-5}.
\end{thebibliography}
\end{document}
